\section{Periodicity and the numerical central charge}
\label{sec:numerical}

Throughout, $X\subset\mathbf P^4_{\C}$ is a smooth hypersurface of degree
five, $i\colon X\hookrightarrow\mathbf P^4$ is its inclusion, and
$\D=\Db(\Coh(X))$. Put $H=c_1(\OO_X(1))$, so that $\int_XH^3=5$.
For an object $E\in\D$ we use the normalized characters
\begin{equation}
\label{eq:normalized-characters}
 v(E)=(r,c,d,e)
 =\left(\rk(E),\frac{\int_XH^2\ch_1(E)}5,
                  \frac{\int_XH\ch_2(E)}5,
                  \frac{\int_X\ch_3(E)}5\right).
\end{equation}
In particular, ordinary slope is normalized as $\mu_H(E)=c(E)/r(E)$
for a sheaf of positive rank. We first record the functor relation that
will make the subsequent family of aisles periodic, and then determine
the central charge on the full numerical Grothendieck group.
The periodicity is already visible in Aspinwall's direct kernel calculation
\cite[Section~7.1.4]{Aspinwall2004}. It also follows from Orlov's equivalence
\cite[Theorem~2.5]{Orlov2009} and its compatibility with grade shift
\cite[Proposition~5.8 and Remark~5.12]{BallardFaveroKatzarkov2012},
as recalled in \cite[Theorem~2.5]{TodaQuintic2013}.
Appendix~\ref{app:periodicity} gives a direct proof, keeping track of the
relative kernels and the natural transformations needed for a functor
isomorphism.  Here we record the relation and compute its numerical action.

\begin{lemma}\label{lem:periodicity}
Let $\Phi=T_{\OO_X}\circ(-\otimes\OO_X(1))$. There is a natural
isomorphism of exact functors
\begin{equation}\label{eq:periodicity}
 \Phi^5\simeq[2].
\end{equation}
Consequently $\Phi$ is an autoequivalence, with inverse $\Phi^4[-2]$.
\end{lemma}

To prove the relation, expand the ordered composition as
\[
 \Phi^5\simeq
 T_{\OO_X}\circ T_{\OO_X(1)}\circ\cdots\circ T_{\OO_X(4)}
 \circ(-\otimes\OO_X(5)).
\]
The rightmost functor acts first.  The relative-kernel argument in
Appendix~\ref{app:periodicity} identifies this twist product with
tensoring by $\OO_X(-5)$ followed by $[2]$.

We next identify the entire numerical group and its integral discreteness.
After computing the charge, these coordinates will also show that controlling
rank, first Chern character, and charge suffices to control a numerical class.

\begin{lemma}\label{lem:numerical}
The map $v$ identifies $\Knum(X)_{\R}$ with $\R^4$. Under this
identification, $\Knum(X)$ is a full discrete lattice contained in
\begin{equation}\label{eq:numerical-lattice}
 \Z\times\Z\times\tfrac1{10}\Z\times\tfrac1{30}\Z.
\end{equation}
For $v=(r,c,d,e)$ and $w=(r',c',d',e')$, the Euler form is
\begin{equation}\label{eq:euler-form}
 \chi(v,w)=5\left(re'-cd'+dc'-er'
                   +\frac56(rc'-cr')\right).
\end{equation}
The matrix induced by $\Phi$ is
\begin{equation}\label{eq:numerical-rotation}
 M=
 \begin{pmatrix}
 -4&-20/3&-5&-5\\
 1&1&0&0\\
 1/2&1&1&0\\
 1/6&1/2&1&1
 \end{pmatrix},
 \qquad
 \det(zI-M)=z^4+z^3+z^2+z+1.
\end{equation}
\end{lemma}

\begin{proof}
The Weak Lefschetz Theorem (in the integral form recalled in
\cite[Theorem~3]{Catanese2014}) and Poincar\'e duality give
$H^{2j}(X,\Q)=\Q H^j$ for $0\le j\le3$; the integral version also
gives $H^2(X,\Z)=\Z H$. Thus the even Chern character is precisely
\[
 \ch(E)=r+cH+dH^2+eH^3.
\]
The tangent sequence gives
$c(T_X)=(1+H)^5/(1+5H)$, so that $c_1(T_X)=0$,
$c_2(T_X)=10H^2$, and
\[
 \operatorname{td}(X)=1+\frac56H^2.
\]
Hirzebruch--Riemann--Roch, obtained by taking the proper map to a point in
\cite[Section~7]{BorelSerre1958}, now gives \eqref{eq:euler-form}, and in
particular
\begin{equation}\label{eq:euler-characteristic}
 \chi(E)=5\left(e(E)+\frac56c(E)\right).
\end{equation}
The matrix of the bilinear form~\eqref{eq:euler-form} is
\[
 J=\begin{pmatrix}
 0&25/6&0&5\\
 -25/6&0&-5&0\\
 0&5&0&0\\
 -5&0&0&0
 \end{pmatrix},
 \qquad\det J=625.
\]
Moreover the four vectors
\[
 v(\OO_X(k))=(1,k,k^2/2,k^3/6),\qquad 0\le k\le3,
\]
span $\Q^4$: their row matrix has determinant one. It follows that
the kernel of $v$ on $K_0(X)$ is exactly the radical of the Euler
pairing. Indeed one implication follows from Riemann--Roch, and the
other follows by testing against these four spanning vectors and using
the nondegeneracy of $J$. Hence $v$ identifies the numerical group with
its image, which spans $\Q^4$.

For any class of $K_0(X)$, integrality of Chern classes and
\[
 \ch_2=\frac{c_1^2-2c_2}{2},\qquad
 \ch_3=\frac{c_1^3-3c_1c_2+3c_3}{6}
\]
give $d\in\tfrac1{10}\Z$ and $e\in\tfrac1{30}\Z$; integral Weak
Lefschetz gives $c\in\Z$. The image is therefore a subgroup of the
discrete lattice in \eqref{eq:numerical-lattice}. Since it contains the
four spanning line-bundle classes, it is itself a full lattice.
In particular its real span is the full space $\R^4$.

Tensoring by $\OO_X(1)$ sends
\[
 (r,c,d,e)\longmapsto
 \left(r,c+r,d+c+\frac r2,e+d+\frac c2+\frac r6\right).
\]
The twist $T_{\OO_X}$ then subtracts $\chi(E(1))$ from the rank and leaves
the other three coordinates unchanged. Formula~\eqref{eq:euler-characteristic}
therefore gives $M$ as displayed. Expanding $\det(zI-M)$ gives
\eqref{eq:numerical-rotation}; in particular its four eigenvalues are
the nonidentity fifth roots of unity, each with multiplicity one.
\end{proof}

We denote by $\Phi_*$ the induced automorphism of $\Knum(X)$ and its
real-linear extension to $\Knum(X)_{\R}$; both are represented by $M$ in
these coordinates.

Put
\begin{equation}\label{eq:charge-constants}
 \lambda=\exp(2\pi i/5),\qquad
 t_0=\frac12\cot(\pi/5)>0,\qquad
 u=\frac45\sin^2(\pi/5)=\frac{5-\sqrt5}{10}.
\end{equation}

\begin{proposition}\label{prop:central-charge}
There is a unique numerical homomorphism $Z\colon\Knum(X)\to\C$
such that $Z\circ\Phi=\lambda Z$ and $Z(\OO_x)=-1$ for closed points
$x\in X$. It is
\begin{equation}\label{eq:central-charge}
 \frac{Z(r,c,d,e)}5
  =\frac{\lambda-1}{5}r
   +\left(t_0^2-\frac56+\frac{i t_0}{2}\right)c
   +\left(-\frac12+i t_0\right)d-e.
\end{equation}
In particular,
\begin{equation}\label{eq:imaginary-charge}
 \frac{\Im Z(r,c,d,e)}{5t_0}=d+\frac12c+ur.
\end{equation}
For any fixed norm on $\Knum(X)_{\R}$, there is a constant $C_0>0$
such that every real numerical class $v=(r,c,d,e)$ satisfies
\begin{equation}\label{eq:numerical-norm-control}
 \norm{v}\le C_0\bigl(|r|+|c|+|Z(v)|\bigr).
\end{equation}
\end{proposition}

\begin{proof}
A closed point has $v(\OO_x)=(0,0,0,1/5)$, so the normalization fixes
the coefficient of $e$ in $Z/5$ to be $-1$.
Write $Z/5=Ar+Bc+Cd-e$. Reading the $e$, $d$, and $c$ columns of
$ZM=\lambda Z$, in that order, gives
\[
 A=\frac{\lambda-1}{5},\qquad
 C=-\frac{\lambda}{\lambda-1},\qquad
 B=-\frac43+\frac{C-1/2}{\lambda-1}.
\]
The identity
\[
 \frac1{\lambda-1}=-\frac12-i t_0
\]
then gives $C=-1/2+i t_0$ and
$B=t_0^2-5/6+i t_0/2$. Substitution in the remaining, rank-column
equation reduces it to
$1+\lambda+\lambda^2+\lambda^3+\lambda^4=0$.
This proves both the eigenvalue identity and uniqueness with the stated
normalization. Taking imaginary parts and using
\[
 \frac{\sin(2\pi/5)}{5t_0}
       =\frac45\sin^2(\pi/5)=u
\]
proves \eqref{eq:imaginary-charge}.

Finally, the real coefficient matrix of $(d,e)$ in
$(\Re Z,\Im Z)$ is
\[
 \begin{pmatrix}-5/2&-5\\5t_0&0\end{pmatrix},
 \qquad \det=25t_0\ne0.
\]
Thus $v\mapsto(r,c,\Re Z(v),\Im Z(v))$ is an isomorphism of real
vector spaces. Boundedness of its inverse in finite dimension proves
\eqref{eq:numerical-norm-control}.
\end{proof}
